% Section 7.3 — Transition to the Paper

\section{Transition to the Paper}
\label{sec:transition_catalog}

The preceding sections laid out the motivation and the conceptual framework: the limitations of fixed-threshold cataloging (Section~\ref{sec:threshold_limits}) and the idea of using the $dN/dS$ as a generative model for simulated-sky comparison (Section~\ref{sec:pop_to_spatial}).
The remainder of this chapter presents the full analysis.

The data selection and background modeling are described in Section~\ref{sec:catalog:data}. The frequentist framework introduced conceptually in Section~\ref{sec:pop_to_spatial} --- the Pearson-like TS, the KS two-sample test, and the Quality Factor --- is formalized in Section~\ref{sec:procedure} and applied to 14 years of \textit{Fermi}-LAT data at high Galactic latitudes ($|b| > 30°$). Five thousand synthetic skies are generated by drawing source fluxes from the $dN/dS$ recovered in Chapter~\ref{ch:6}, and their pixel-wise TS distributions are compared against the real sky via two-sample Kolmogorov--Smirnov tests. The results, presented in Section~\ref{sec:catalogue}, explore the full $(\alpha, \mathrm{QF})$ plane, providing firing pixel counts for multiple operating points.
For the conventional choice $\alpha = 0.05$ and a quality factor $\mathrm{QF} = 0.5$, the procedure identifies ${\sim}\,9{,}600$ firing pixels, compared to ${\sim}\,6{,}700$ inferred from the 4FGL-DR3 catalog map alone --- an increase of roughly 50\%.
As a sanity check, the fraction of bright 4FGL-DR3 sources recovered as firing pixels approaches unity for fluxes above ${\sim}\,2 \times 10^{-10}\;\mathrm{cm^{-2}\,s^{-1}}$, where \textit{Fermi}-LAT's standard detection sensitivity begins to saturate.
The robustness of the results against the choice of Galactic foreground model is also assessed, with pixel-level agreement exceeding 97\% across independent background prescriptions (Section~\ref{sec:catalogue}).

The results are publicly available as the \texttt{gPCS} (\textit{gamma-ray Photon Count Statistics}) Python package, together with a summary FITS file containing pixel indices, TS values, and quality factors.
Both are hosted on Zenodo and GitHub, allowing users to select the $(\alpha, \mathrm{QF})$ operating point most appropriate to a given application.
One natural use of this extended catalog is cross-correlation studies between gamma-ray maps and tracers of large-scale structure, where the statistical advantage of a larger sample of candidate directions outweighs a sub-leading fraction of spurious entries --- the subject of Chapter~\ref{ch:8}. Future directions, including energy-dependent extensions and low-latitude applications, are discussed in Section~\ref{sec:catalog:conclusions}.
